% =================================================================
\section{Preguntas de análisis}
% =================================================================
\paragraph{1. ¿Cuántos servidores para 300 peticiones/s con respuestas en menos de medio segundo?}
Un servidor atiende 112 peticiones/s en menos de medio segundo, pero en ese punto su CPU está
casi al 100\,\%. Para tener margen conviene trabajar al 70\,\% (unas 62 peticiones/s por
servidor), lo que da $\lceil 300/62\rceil = $ \textbf{5 servidores}. Eso consume exactamente los
10 núcleos de nuestra cuota, así que \textbf{es viable, pero sin margen}. Para soportar además la
caída de una zona harían falta 6 (dos por zona). Lo que cambiaríamos: pedir más cuota a Azure
(al menos 16 núcleos), subir el máximo del autoscale a 6--8 y hacer la aplicación más eficiente
(caché de páginas, CDN para imágenes y un servidor web más rápido).

\paragraph{2. ¿Cuánto tardó el autoscale y qué hacer ante un pico a las 8:00 p.\,m.?}
Desde que la CPU superó el 70\,\% pasaron \textbf{5 min 19 s hasta que se ordenó el nuevo servidor}
y unos 6 minutos hasta que atendió tráfico. Ese tiempo es la suma de la ventana de 5 minutos que
usa para decidir, la demora de la métrica, la creación de la máquina y su instalación. Ante un pico
repentino serían varios minutos de lentitud. Proponemos:
\begin{itemize}
  \item \textbf{Perfil programado:} de 7:30 a 10:30 p.\,m. mantener al menos 4 servidores, para que
        la capacidad ya esté lista cuando empiece la campaña.
  \item \textbf{Autoscale predictivo:} Azure aprende de los últimos 7 días y agrega servidores antes
        del pico. Primero lo usaríamos solo en modo pronóstico para validarlo.
  \item \textbf{Reglas más rápidas} solo en ese horario: decidir con 2 minutos de datos, umbral del
        60\,\% y agregar 2 servidores a la vez.
\end{itemize}

\paragraph{3. ¿Cuál es ahora el punto único de falla?}
Dentro de la región no hay uno: la IP, el balanceador y los servidores están repartidos en tres
zonas. El punto único de falla es \textbf{la propia región East US}, porque todo vive en ella, y en
menor medida \textbf{el Key Vault}, que es uno solo (la aplicación guarda la contraseña 60\,s en
memoria, lo que solo cubre fallas cortas). Eliminarlo implica una copia de la plataforma en otra
región con Azure Front Door delante: unos USD 35/mes de Front Door más la segunda plataforma
(USD 100 a 174/mes). En total, entre 135 y 210 dólares más al mes, casi duplicar el costo.

\paragraph{4. Si un atacante lograra ejecutar código en un servidor, ¿qué podría hacer?}
\textbf{Podría} leer la contraseña de la base de datos (es la única que ese servidor puede leer),
alterar las respuestas de ese servidor o consumir su CPU. \textbf{No podría}:
\begin{itemize}
  \item leer otros secretos: el permiso está limitado a un solo secreto;
  \item cambiar la infraestructura: el servidor no tiene permisos de administración en Azure;
  \item entrar a otros servidores por SSH: no hay IPs públicas y la red solo deja pasar la web;
  \item enviar datos a un servidor propio en Internet: la salida está bloqueada (prueba P3.2);
  \item usar el almacenamiento: no tiene permisos y las llaves de acceso están desactivadas.
\end{itemize}
\textbf{Riesgo que queda:} la salida permitida hacia Storage abarca cualquier cuenta de la región,
incluso una del atacante. Se cerraría con conexiones privadas o un firewall que filtre por dominio.
Además, cada lectura del secreto queda registrada (S6), y tras un incidente se puede cambiar la
contraseña con un solo comando de Terraform.

\paragraph{5. ¿Por qué Checkov en el pipeline no reemplaza a Azure Policy, ni al revés?}
Porque protegen momentos distintos. \textbf{Checkov} revisa el código \emph{antes} de desplegar y
avisa temprano, pero solo ve lo que pasa por el repositorio: un cambio hecho a mano en el portal no
lo atraviesa, y un \texttt{\#checkov:skip} puede silenciarlo. \textbf{Azure Policy} actúa
\emph{dentro de Azure} y bloquea a cualquiera, venga de Terraform, del portal o de un atacante, pero
solo avisa en el momento de crear el recurso. Juntas cubren ambos frentes: en P3.3 la política
rechazó recursos creados por fuera de Terraform.

\paragraph{6. El estado de Terraform contiene secretos: ¿cómo se protegió?}
El archivo de estado guarda la contraseña generada, la llave SSH y el mapa de la infraestructura.
Lo protegimos así:
\begin{itemize}
  \item vive en Azure, nunca en el computador ni en Git;
  \item la cuenta no acepta llaves de acceso: solo entra quien tenga identidad y permiso en Azure;
  \item firewall limitado a nuestras IPs, cifrado y conexión segura;
  \item historial de versiones y papelera de 7 días, y bloqueo para que dos personas no lo editen
        a la vez.
\end{itemize}
Si se filtrara, el atacante tendría la contraseña y la llave SSH (inútil, porque el puerto 22 no es
alcanzable) y conocería la arquitectura. Habría que cambiar la contraseña, regenerar la llave y
revisar los registros de acceso. Como mejora, las versiones recientes de Terraform permiten que
los secretos nunca se guarden en el estado.

% =================================================================
\section{Conclusiones}
% =================================================================
La plataforma resuelve los tres problemas de CaribeShop, y cada solución se comprobó con datos.
Ante la \textbf{saturación}, el autoscale pasó de 2 a 4 servidores y volvió a 2 con solo 0,55\,\%
de errores. Ante la \textbf{falla de una máquina}, un servidor enfermo salió de circulación en
8 segundos sin que los usuarios lo notaran, y cuando borramos uno de golpe la capacidad se
recuperó en menos de 2 minutos. Ante la \textbf{auditoría}, ya no hay secretos en el código ni
accesos de administración expuestos, y el pipeline bloqueó un cambio que intentaba abrir SSH.

La seguridad se construyó en capas, y cada una se puso a prueba intentando romperla. Medir cuánto
aguanta un servidor nos permitió justificar la configuración del autoscale y responder cuántos se
necesitan para 300 peticiones por segundo. La principal limitación es que el autoscale tarda unos
5 a 6 minutos en reaccionar, algo que en una campaña se resuelve programándolo con anticipación.
Al terminar eliminamos todos los recursos con \texttt{terraform destroy}. Las capturas y salidas
de comandos que respaldan cada requisito están en la carpeta \texttt{evidencias/} del repositorio.
